% =====================================================================
%  ANEXOS
%  Anexo A, añadido el 27 de octubre de 2026 a pedido de los autores: los
%  videos del proyecto publicados en el canal de YouTube de Santiago
%  Hernández, en orden de publicación. Están «no listados» (se abren con
%  el enlace) salvo el primero, que es público. Los otros videos del canal
%  son de otras asignaturas y no entran. La descripción de cada uno sigue la
%  de su entregable semanal o la del sitio del proyecto.
%  Anexo B, añadido el mismo día: el repositorio, con toda la trazabilidad.
% =====================================================================
\chapter*{Anexo A. Videos del proyecto}
\addcontentsline{toc}{chapter}{ANEXO A. VIDEOS DEL PROYECTO}

Videos publicados en YouTube durante el proyecto, en orden de publicación. Cada
código QR abre el video correspondiente.

\begin{table}[H]
  \centering\small
  \renewcommand{\arraystretch}{1.05}
  \begin{tabular}{@{}c >{\raggedright\arraybackslash}p{0.13\textwidth} p{0.52\textwidth} c c@{}}
    \toprule
    \textbf{N.\textsuperscript{o}} & \textbf{Publicado} & \textbf{Contenido} & \textbf{Duración} & \textbf{Enlace} \\ \midrule
    1 & 5 may. 2026 & El agente en el modelo simulado del pasillo de la universidad, con el seguidor de paredes. \newline \url{https://youtu.be/Q0w_CTga98k} & 1:12 & \raisebox{-0.85\height}{\qrcode[height=1.25cm]{https://youtu.be/Q0w_CTga98k}} \\ \addlinespace
    2 & 19 may. 2026 & Teleoperación del agente en el modelo del piso 2 del edificio Gregorio, con la respuesta de los sensores y el mapeo en RViz. \newline \url{https://youtu.be/ZPJ0Qw2L5-A} & 1:00 & \raisebox{-0.85\height}{\qrcode[height=1.25cm]{https://youtu.be/ZPJ0Qw2L5-A}} \\ \addlinespace
    3 & 5 ago. 2026 & Liberación y prueba de marcha del AWS DeepRacer desde su interfaz web. \newline \url{https://youtu.be/ZGfAMnC4lYY} & 0:44 & \raisebox{-0.85\height}{\qrcode[height=1.25cm]{https://youtu.be/ZGfAMnC4lYY}} \\ \addlinespace
    4 & 18 ago. 2026 & Media vuelta en simulación con Nav2: la meta queda detrás del agente y el planificador la resuelve con una inversión de marcha, vista en Gazebo y en RViz. \newline \url{https://youtu.be/6SABSgCxVEw} & 0:47 & \raisebox{-0.85\height}{\qrcode[height=1.25cm]{https://youtu.be/6SABSgCxVEw}} \\ \addlinespace
    5 & 24 ago. 2026 & Primeras asignaciones de destinos desde el coordinador, con un agente y Nav2 en simulación. \newline \url{https://youtu.be/iKYSiM9cs-Y} & 0:45 & \raisebox{-0.85\height}{\qrcode[height=1.25cm]{https://youtu.be/iKYSiM9cs-Y}} \\ \addlinespace
    6 & 5 sep. 2026 & Campaña de treinta misiones sorteadas en simulación, con los dos agentes y el coordinador. \newline \url{https://youtu.be/eH_Co4vZZdw} & 3:48 & \raisebox{-0.85\height}{\qrcode[height=1.25cm]{https://youtu.be/eH_Co4vZZdw}} \\ \addlinespace
    7 & 20 oct. 2026 & Campaña con los agentes físicos en los pisos 3 y 4 del edificio Catalayud: misiones con relevo pedidas, seguidas y confirmadas desde un teléfono. \newline \url{https://youtu.be/X3BINGgNxiA} & 12:40 & \raisebox{-0.85\height}{\qrcode[height=1.25cm]{https://youtu.be/X3BINGgNxiA}} \\ \bottomrule
  \end{tabular}
\end{table}

\chapter*{Anexo B. Repositorio del proyecto}
\addcontentsline{toc}{chapter}{ANEXO B. REPOSITORIO DEL PROYECTO}

El repositorio público del proyecto guarda toda su trazabilidad. Contiene el
código de los agentes, del coordinador y de la interfaz; la matriz de requisitos,
con la prueba que declara cumplido cada uno; el protocolo experimental; el
registro de cada misión de las dos campañas y el programa que calcula las
métricas; la evidencia de cada sesión de trabajo, con sus mediciones; los
entregables semanales y la bitácora del proyecto. El sitio del proyecto presenta
el mismo contenido para leerlo en el navegador.

\begin{table}[H]
  \centering\small
  \renewcommand{\arraystretch}{1.3}
  \begin{tabular}{@{}p{0.2\textwidth} p{0.56\textwidth} c@{}}
    \toprule
    \textbf{Recurso} & \textbf{Dirección} & \textbf{Código} \\ \midrule
    Repositorio & \url{https://github.com/shdeza12/Collaborative-Mobile-Robot-System-for-Indoor-Navigation-Assistance-in-Multi-Level-Environments} & \raisebox{-0.85\height}{\qrcode[height=2cm]{https://github.com/shdeza12/Collaborative-Mobile-Robot-System-for-Indoor-Navigation-Assistance-in-Multi-Level-Environments}} \\ \addlinespace
    Sitio del proyecto & \url{https://shdeza12.github.io/Collaborative-Mobile-Robot-System-for-Indoor-Navigation-Assistance-in-Multi-Level-Environments/} & \raisebox{-0.85\height}{\qrcode[height=2cm]{https://shdeza12.github.io/Collaborative-Mobile-Robot-System-for-Indoor-Navigation-Assistance-in-Multi-Level-Environments/}} \\ \bottomrule
  \end{tabular}
\end{table}
